% Generated by Listas/make_statements.py from pratica-lecture-01.md, item 5. Do not edit.
\textbf{Source.} Muñoz-Garcia \& Toro-Gonzalez (2019), \emph{Strategy and Game Theory: Practice Exercises with Answers}, Ch. 1, Exercise 13 ``Iterated Deletion of Strictly Dominated Strategies (IDSDS)'', p.~26 {[}PDF 39{]}. Reproduced verbatim (statement only; the book's worked answer follows it).



\begin{quote}

Consider the following normal form game (Fig. 1.33)



{\def\LTcaptype{none} % do not increment counter

\begin{longtable}[]{@{}lll@{}}

\toprule\noalign{}

Player 1 ~Player 2 & \emph{Left} & \emph{Right} \\

\midrule\noalign{}

\endhead

\bottomrule\noalign{}

\endlastfoot

\emph{Up} & -10, 10 & 0, 12 \\

\emph{Middle} & -12, 0 & 2, 2 \\

\emph{Down} & -9, 0 & 1, 0 \\

\end{longtable}

}



\emph{Fig. 1.33} Normal-form game



\begin{enumerate}

\def\labelenumi{(\alph{enumi})}

\item

  Find strictly dominant strategies (if any) for player 1 and for player 2.

\item

  Find strictly dominated strategies (if any) for player 1 and for player 2.

\item

  If you apply iterative deletion of strictly dominated strategies (IDSDS), what is the surviving strategy pair (or pairs)? Explain the steps you use in IDSDS, and why you use them.

\end{enumerate}

\end{quote}

